% Source: Wald GR, physical PDF pp. 391--393
%         (printed pp. 378--380).
% Coverage: Chapter 14 opening material before Section 14.1.
% Figures/tables/footnotes: no figures, tables, or footnotes.
% Status: source-reviewed and compile-clean.
% Uncertainties: none.

\setcounter{footnote}{0}

As described in chapter~4, the theory of general relativity put forth a revolutionary
new viewpoint on spacetime structure and gravitation. However, in one important
sense, this theory is not revolutionary enough. It is well established that all known
physical fields must be described on a fundamental level by the principles of quantum
theory. In quantum theory, states of a system are represented by vectors in a Hilbert
space, \(\mathcal H\), and observable quantities are represented by self-adjoint linear maps
acting on \(\mathcal H\). Unless the state of the system happens to be in an eigenstate of the
observable, the observable will not have a definite value and one can predict only
probabilities for the outcomes of measurements. However, general relativity is a
purely classical theory, since in the framework of general relativity the observable
quantities---in particular, the spacetime metric---always have definite values. Thus,
if the principles of quantum theory are to apply to the gravitational field, general
relativity must at best be only an approximation to a truly fundamental theory of
gravity, perhaps in the same sort of way as Maxwell's theory of electromagnetism
is only an approximation to quantum electrodynamics.

Classical descriptions of ordinary matter are, in general, excellent approximations
for describing phenomena which occur on macroscopic scales, although even here
quantum effects can be important for suitably prepared states. However, the classical
description of matter becomes wholly inadequate on atomic and smaller scales. In
this case, the scale at which the classical description breaks down is determined by
the masses and charges of the fundamental particles as well as the two fundamental
constants of nature which enter the theory, namely Planck's constant, \(\hbar\), and the
speed of light, \(c\). Similarly, in a quantum theory of gravitation based on general
relativity, one would expect that the fundamental scale at which the classical description
becomes wholly inadequate should be set by \(\hbar\), \(c\), and the gravitational constant,
\(G\). There is a unique combination of these constants which has the dimensions of
length, namely the quantity
\(\ell_P\equiv(G\hbar/c^3)^{1/2}\), called the \emph{Planck length}. As might be
expected, the Planck length arises naturally in attempts to formulate a quantum
theory of gravity. Thus, dimensional arguments suggest that a classical description
of spacetime structure should break down at scales of the order of the Planck length
and smaller.

In cgs units the magnitude of the Planck length is only \(\sim10^{-33}\,\mathrm{cm}\).
(The corresponding Planck scales of other quantities such as time and energy are given in
appendix~F.) The smallness of \(\ell_P\) compared with typical length scales occurring in
atomic, nuclear, and elementary particle physics is directly related to the weakness
of the gravitational force between two elementary particles as compared with the
other fundamental forces, namely the strong, weak, and electromagnetic interactions.
The Planck scales lie many orders of magnitude beyond what we presently
are able to probe with high energy accelerators. Therefore, it might appear that the
development of a quantum theory of gravity---while undoubtedly a laudable goal---
would be unlikely to have much relevance to any presently observable phenomena.

However, there are at least two good reasons for believing that the predictions of
a quantum theory of gravity could be very relevant to phenomena at presently
observable scales. The first reason arises from the theory of elementary particles. It
is widely believed that the true, fundamental theory of nature will achieve a
unification of the forces of nature by describing them simply as different aspects of
a single entity. A unification of the classical theories of electricity and magnetism
was achieved over a century ago by Maxwell. More recently, the Weinberg--Salam
theory has given a successful unified description of the weak and electromagnetic
interactions. It is presently believed that the ``grand unified models'' may successfully
unify the strong and electroweak interactions. The unification of gravitation with the
strong--electroweak interaction would be the next logical step in this program. Interestingly,
the natural length scale which arises in the grand unified theories is only a
few orders of magnitude larger than \(\ell_P\). Thus, it is quite possible that a quantum
theory of gravitation may even play an important role in the unification of the strong
and electroweak interactions. A unified theory of all forces undoubtedly would yield
many new predictions of phenomena at presently observable scales.

The second reason arises directly from general relativity. As discussed in chapter~9,
spacetime singularities occur in the solutions of classical general relativity relevant
to gravitational collapse and cosmology. Thus, in these situations, the classical
description of spacetime structure must break down. In particular, one cannot expect
the homogeneous, isotropic models of chapter~5 to be an adequate description of our
universe in the regime where they predict curvature of magnitude \(\ell_P^{-2}\) or greater,
i.e., for
\(t<t_P=(G\hbar/c^5)^{1/2}\sim10^{-43}\,\mathrm{s}\). Thus, it appears that the
development of a quantum theory of gravitation will be an essential requirement for
our understanding of the initial state of our universe. It is not implausible that phenomena
which occur in the very early universe and which can be understood only in the framework of
quantum gravity will lead to observationally verifiable predictions about the structure
of the present universe.

However, even if quantum gravity leads to no predictions of phenomena which
can be observed with present day technology, the formulation of a quantum theory
of gravity undoubtedly would be of major significance for theoretical physics. After
all, classical general relativity has provided us with major new insights into the
workings of nature even though the new observationally verifiable predictions it
makes are relatively meager. The new fundamental insights provided by a quantum
theory of gravity certainly should be no less significant than those provided by
general relativity.

As discussed in section~14.1, all of the known procedures for formulating a
quantum field theory associated with a classical theory run into difficulties when
applied to general relativity. Thus, the formulation of a viable quantum theory of
gravity remains a goal for the future. However, a completely satisfactory theory
exists for a free (i.e., linear) quantum matter field propagating in a fixed background
curved spacetime. Although such a theory is, at best, only an approximation to a full
quantum theory of gravity with quantum matter, the effects thereby predicted should
at least give a good indication of the types of quantum effects which may occur in
strong gravitational fields. In particular, as described in section~14.2, the creation of
particles by a gravitational field is predicted by this theory. Remarkably, when
applied to the case of a black hole in section~14.3, one finds that particle creation
causes the effective ``emission'' by a black hole of a thermal spectrum of particles at
temperature \(kT=\hbar\kappa/(2\pi)\), where \(\kappa\) is the surface gravity of the black hole.
The implications of this result for the relationship between black holes and
thermodynamics are explored in section~14.4.
